
\documentclass{article}
\usepackage{colortbl}
\usepackage{makecell}
\usepackage{multirow}
\usepackage{supertabular}

\begin{document}

\newcounter{utterance}

\centering \large Interaction Transcript for game `dond', experiment `coop\_en', episode 5 with reflexion{-}llama{-}3.1{-}8b{-}instruct.
\vspace{24pt}

{ \footnotesize  \setcounter{utterance}{1}
\setlength{\tabcolsep}{0pt}
\begin{supertabular}{c@{$\;$}|p{.15\linewidth}@{}p{.15\linewidth}p{.15\linewidth}p{.15\linewidth}p{.15\linewidth}p{.15\linewidth}}
    \# & $\;$A & \multicolumn{4}{c}{Game Master} & $\;\:$B\\
    \hline

    \theutterance \stepcounter{utterance}  
    & & \multicolumn{4}{p{0.6\linewidth}}{
        \cellcolor[rgb]{0.9,0.9,0.9}{
            \makecell[{{p{\linewidth}}}]{
                \texttt{\tiny{[P1$\langle$GM]}}
                \texttt{You have attempted this same task before and did not succeed. These are the plans you wrote after those attempts.} \\
\texttt{Trial \#1: New plan: } \\
\\ 
\texttt{In my previous attempt, I failed to achieve a mutually beneficial deal because I didn't effectively communicate my priorities and didn't consider the other player's perspective. I also didn't use my messages efficiently, as I spent too much time discussing the cups and hats, which were not as valuable to me. } \\
\\ 
\texttt{Next time, I will start by clearly stating my priorities and the items I am willing to compromise on. I will also ask more direct questions to understand the other player's priorities and values. Additionally, I will try to find a mutually beneficial deal earlier in the conversation, rather than waiting until the end. } \\
\\ 
\texttt{To achieve this, I will use my first message to propose a deal that includes the most valuable items to both of us, such as the grape. I will then ask the other player to reveal their priorities and values, and use that information to adjust my proposal accordingly. I will also try to find a way to compromise on the less valuable items, such as the hats, to make the deal more mutually beneficial.} \\
\\ 
\texttt{Trial \#2: New plan: } \\
\\ 
\texttt{In my previous attempt, I failed to achieve a mutually beneficial deal because I didn't effectively communicate my priorities and didn't consider the other player's perspective. I also didn't use my messages efficiently, as I spent too much time discussing the cups and hats, which were not as valuable to me. } \\
\\ 
\texttt{However, I noticed that I was able to reach a mutually beneficial deal when I proposed a slight modification to the other player's offer. This suggests that I should focus on finding areas of compromise and adjusting my proposals accordingly.} \\
\\ 
\texttt{Next time, I will start by clearly stating my priorities and the items I am willing to compromise on. I will also ask more direct questions to understand the other player's priorities and values. } \\
\\ 
\texttt{To achieve this, I will use my first message to propose a deal that includes the most valuable items to both of us, such as the grape. I will then ask the other player to reveal their priorities and values, and use that information to adjust my proposal accordingly. } \\
\\ 
\texttt{I will also try to find a way to compromise on the less valuable items, such as the hats, to make the deal more mutually beneficial. Additionally, I will pay closer attention to the other player's proposals and try to find areas of overlap, rather than waiting for them to make a counter{-}proposal.} \\
\\ 
\\ 
\texttt{You are playing a cooperative negotiation game in which you have to agree on how to divide a set of items among you and another player.} \\
\\ 
\texttt{Rules:} \\
\texttt{(a) You and the other player are given a set of items. Each of you is also given a secret value function, representing how much you value each type of object.} \\
\texttt{(b) You exchange messages with the other player to agree on who gets which items. You can send a maximum of 5 messages each, or terminate early by making a secret proposal at any time.} \\
\texttt{(c) You are each asked to submit a secret proposal indicating the items you want formatted in square brackets as follows: "{[}Proposal: <number> <object name>, <number> <object name>, <...>{]}"} \\
\texttt{(d) If your proposals are complementary, i.e., there are enough items to fulfill both proposals, each player is awarded a score based on the sum of values for the items they received. Otherwise, both of you get zero points.} \\
\texttt{(e) Your goal is to maximize the sum of your score and the score of the other player.} \\
\\ 
\texttt{Let us start.} \\
\\ 
\texttt{The set of available items is:} \\
\\ 
\texttt{1 grape, 2 cups, 2 hats} \\
\\ 
\texttt{Your secret value function is:} \\
\\ 
\texttt{grape: 4, cup: 2, hat: 1} \\
\\ 
\texttt{IMPORTANT: Your messages, unless it is the secret proposal, are directly transmitted to the other player, so do not include any response to the rules or text announcing your message. To make a secret proposal, use the indicated format. Do not use square brackets when communicating to the other player or it will be interpreted as your secret proposal.} \\
            }
        }
    }
    & \\ \\

    \theutterance \stepcounter{utterance}  
    & \multicolumn{4}{p{0.6\linewidth}}{
        \cellcolor[rgb]{0.9,0.9,0.9}{
            \makecell[{{p{\linewidth}}}]{
                \texttt{\tiny{[P1$\rangle$GM]}}
                \texttt{I'd like to propose a deal that includes the most valuable items to both of us. I think the grape is a good starting point, as it has the highest value for both of us. How about we start with a proposal that includes the grape and one of the cups? This way, we can both get something we value highly, and then we can discuss the remaining items. } \\
\\ 
\texttt{What do you think about this proposal: 1 grape, 1 cup?} \\
            }
        }
    }
    & & \\ \\

    \theutterance \stepcounter{utterance}  
    & & & \multicolumn{2}{p{0.3\linewidth}}{
        \cellcolor[rgb]{0.9,0.9,0.9}{
            \makecell[{{p{\linewidth}}}]{
                \texttt{\tiny{[GM$|$GM]}}
                \texttt{continue} \\
            }
        }
    }
    & & \\ \\

    \theutterance \stepcounter{utterance}  
    & & & \multicolumn{2}{p{0.3\linewidth}}{
        \cellcolor[rgb]{0.9,0.9,0.9}{
            \makecell[{{p{\linewidth}}}]{
                \texttt{\tiny{[GM$|$GM]}}
                \texttt{I'd like to propose a deal that includes the most valuable items to both of us. I think the grape is a good starting point, as it has the highest value for both of us. How about we start with a proposal that includes the grape and one of the cups? This way, we can both get something we value highly, and then we can discuss the remaining items. } \\
\\ 
\texttt{What do you think about this proposal: 1 grape, 1 cup?} \\
            }
        }
    }
    & & \\ \\

    \theutterance \stepcounter{utterance}  
    & & \multicolumn{4}{p{0.6\linewidth}}{
        \cellcolor[rgb]{0.9,0.9,0.9}{
            \makecell[{{p{\linewidth}}}]{
                \texttt{\tiny{[GM$\rangle$P2]}}
                \texttt{You have attempted this same task before and did not succeed. These are the plans you wrote after those attempts.} \\
\texttt{Trial \#1: New plan: } \\
\\ 
\texttt{In my previous attempt, I failed to achieve a mutually beneficial deal because I didn't effectively communicate my priorities and didn't make a clear proposal that would satisfy both parties. } \\
\\ 
\texttt{To improve, I will start by making a clear and concise proposal in my first message, indicating the items I am willing to give up and the items I want to keep. This will give the other player a clear understanding of my priorities and allow us to start negotiating from a common ground.} \\
\\ 
\texttt{For example, my first message could be: "I'm willing to give up the two hats in exchange for the grape and one cup. Is this a deal you would be willing to accept?"} \\
\\ 
\texttt{This approach will allow us to quickly identify areas of agreement and disagreement, and make it easier to find a mutually beneficial solution.} \\
\\ 
\texttt{Trial \#2: New plan: } \\
\\ 
\texttt{In my previous attempt, I failed to achieve a mutually beneficial deal because I didn't effectively communicate my priorities and didn't make a clear proposal that would satisfy both parties. However, I also didn't take into account the other player's priorities and values, which led to a series of back{-}and{-}forth proposals that didn't ultimately result in a mutually beneficial deal.} \\
\\ 
\texttt{To improve, I will start by making a clear and concise proposal in my first message, indicating the items I am willing to give up and the items I want to keep. I will also take into account the other player's priorities and values, and try to find a compromise that satisfies both parties.} \\
\\ 
\texttt{Additionally, I will pay closer attention to the other player's responses and adjust my proposals accordingly. If the other player is willing to give up a certain item, I will try to incorporate that item into my proposal. If the other player is unwilling to give up a certain item, I will try to find an alternative solution that still satisfies both parties.} \\
\\ 
\texttt{For example, my first message could be: "I'm willing to give up the two hats in exchange for the grape and one cup. Is this a deal you would be willing to accept?"} \\
\\ 
\texttt{However, I will also be prepared to adjust my proposal based on the other player's response. If they are willing to give up the grape, I may propose giving up the two cups instead. If they are unwilling to give up the grape, I may propose giving up the two hats and one cup instead.} \\
\\ 
\texttt{By taking a more flexible and responsive approach, I hope to be able to find a mutually beneficial deal that satisfies both parties.} \\
\\ 
\\ 
\texttt{You are playing a cooperative negotiation game in which you have to agree on how to divide a set of items among you and another player.} \\
\\ 
\texttt{Rules:} \\
\texttt{(a) You and the other player are given a set of items. Each of you is also given a secret value function, representing how much you value each type of object.} \\
\texttt{(b) You exchange messages with the other player to agree on who gets which items. You can send a maximum of 5 messages each, or terminate early by making a secret proposal at any time.} \\
\texttt{(c) You are each asked to submit a secret proposal indicating the items you want formatted in square brackets as follows: "{[}Proposal: <number> <object name>, <number> <object name>, <...>{]}"} \\
\texttt{(d) If your proposals are complementary, i.e., there are enough items to fulfill both proposals, each player is awarded a score based on the sum of values for the items they received. Otherwise, both of you get zero points.} \\
\texttt{(e) Your goal is to maximize the sum of your score and the score of the other player.} \\
\\ 
\texttt{Let us start.} \\
\\ 
\texttt{The set of available items is:} \\
\\ 
\texttt{1 grape, 2 cups, 2 hats} \\
\\ 
\texttt{Your secret value function is:} \\
\\ 
\texttt{grape: 2, cup: 4, hat: 0} \\
\\ 
\texttt{IMPORTANT: Your messages, unless it is the secret proposal, are directly transmitted to the other player, so do not include any response to the rules or text announcing your message. To make a secret proposal, use the indicated format. Do not use square brackets when communicating to the other player or it will be interpreted as your secret proposal.} \\
\\ 
\\ 
\texttt{I'd like to propose a deal that includes the most valuable items to both of us. I think the grape is a good starting point, as it has the highest value for both of us. How about we start with a proposal that includes the grape and one of the cups? This way, we can both get something we value highly, and then we can discuss the remaining items. } \\
\\ 
\texttt{What do you think about this proposal: 1 grape, 1 cup?} \\
            }
        }
    }
    & \\ \\

    \theutterance \stepcounter{utterance}  
    & & & \multicolumn{4}{p{0.6\linewidth}}{
        \cellcolor[rgb]{0.9,0.9,0.9}{
            \makecell[{{p{\linewidth}}}]{
                \texttt{\tiny{[GM$\langle$P2]}}
                \texttt{I think the grape is a good starting point, and including one cup in the proposal is a good compromise. However, I'm not sure about giving up one cup, as it has a higher value for me than the grape has for you. } \\
\\ 
\texttt{How about we modify the proposal to include the grape and one hat instead? This way, we can both get something we value, and we can discuss the remaining items. What do you think about this proposal: 1 grape, 1 hat?} \\
            }
        }
    }
     \\ \\

    \theutterance \stepcounter{utterance}  
    & & & \multicolumn{2}{p{0.3\linewidth}}{
        \cellcolor[rgb]{0.9,0.9,0.9}{
            \makecell[{{p{\linewidth}}}]{
                \texttt{\tiny{[GM$|$GM]}}
                \texttt{continue} \\
            }
        }
    }
    & & \\ \\

    \theutterance \stepcounter{utterance}  
    & & & \multicolumn{2}{p{0.3\linewidth}}{
        \cellcolor[rgb]{0.9,0.9,0.9}{
            \makecell[{{p{\linewidth}}}]{
                \texttt{\tiny{[GM$|$GM]}}
                \texttt{I think the grape is a good starting point, and including one cup in the proposal is a good compromise. However, I'm not sure about giving up one cup, as it has a higher value for me than the grape has for you. } \\
\\ 
\texttt{How about we modify the proposal to include the grape and one hat instead? This way, we can both get something we value, and we can discuss the remaining items. What do you think about this proposal: 1 grape, 1 hat?} \\
            }
        }
    }
    & & \\ \\

    \theutterance \stepcounter{utterance}  
    & & \multicolumn{4}{p{0.6\linewidth}}{
        \cellcolor[rgb]{0.9,0.9,0.9}{
            \makecell[{{p{\linewidth}}}]{
                \texttt{\tiny{[P1$\langle$GM]}}
                \texttt{I think the grape is a good starting point, and including one cup in the proposal is a good compromise. However, I'm not sure about giving up one cup, as it has a higher value for me than the grape has for you. } \\
\\ 
\texttt{How about we modify the proposal to include the grape and one hat instead? This way, we can both get something we value, and we can discuss the remaining items. What do you think about this proposal: 1 grape, 1 hat?} \\
            }
        }
    }
    & \\ \\

    \theutterance \stepcounter{utterance}  
    & \multicolumn{4}{p{0.6\linewidth}}{
        \cellcolor[rgb]{0.9,0.9,0.9}{
            \makecell[{{p{\linewidth}}}]{
                \texttt{\tiny{[P1$\rangle$GM]}}
                \texttt{I see your point about the cup having a higher value for you. Including a hat in the proposal is a good compromise, and it's a good way to start building a mutually beneficial deal.} \\
\\ 
\texttt{I think the proposal of 1 grape, 1 hat is a good starting point. However, I'm not sure about giving up the hat, as it has a relatively low value for me. How about we modify the proposal to include the grape and one of the cups, and then we can discuss the remaining items? This way, we can both get something we value highly, and we can find a way to compromise on the remaining items.} \\
\\ 
\texttt{What do you think about this proposal: 1 grape, 1 cup?} \\
            }
        }
    }
    & & \\ \\

    \theutterance \stepcounter{utterance}  
    & & & \multicolumn{2}{p{0.3\linewidth}}{
        \cellcolor[rgb]{0.9,0.9,0.9}{
            \makecell[{{p{\linewidth}}}]{
                \texttt{\tiny{[GM$|$GM]}}
                \texttt{continue} \\
            }
        }
    }
    & & \\ \\

    \theutterance \stepcounter{utterance}  
    & & & \multicolumn{2}{p{0.3\linewidth}}{
        \cellcolor[rgb]{0.9,0.9,0.9}{
            \makecell[{{p{\linewidth}}}]{
                \texttt{\tiny{[GM$|$GM]}}
                \texttt{I see your point about the cup having a higher value for you. Including a hat in the proposal is a good compromise, and it's a good way to start building a mutually beneficial deal.} \\
\\ 
\texttt{I think the proposal of 1 grape, 1 hat is a good starting point. However, I'm not sure about giving up the hat, as it has a relatively low value for me. How about we modify the proposal to include the grape and one of the cups, and then we can discuss the remaining items? This way, we can both get something we value highly, and we can find a way to compromise on the remaining items.} \\
\\ 
\texttt{What do you think about this proposal: 1 grape, 1 cup?} \\
            }
        }
    }
    & & \\ \\

    \theutterance \stepcounter{utterance}  
    & & \multicolumn{4}{p{0.6\linewidth}}{
        \cellcolor[rgb]{0.9,0.9,0.9}{
            \makecell[{{p{\linewidth}}}]{
                \texttt{\tiny{[GM$\rangle$P2]}}
                \texttt{I see your point about the cup having a higher value for you. Including a hat in the proposal is a good compromise, and it's a good way to start building a mutually beneficial deal.} \\
\\ 
\texttt{I think the proposal of 1 grape, 1 hat is a good starting point. However, I'm not sure about giving up the hat, as it has a relatively low value for me. How about we modify the proposal to include the grape and one of the cups, and then we can discuss the remaining items? This way, we can both get something we value highly, and we can find a way to compromise on the remaining items.} \\
\\ 
\texttt{What do you think about this proposal: 1 grape, 1 cup?} \\
            }
        }
    }
    & \\ \\

    \theutterance \stepcounter{utterance}  
    & & & \multicolumn{4}{p{0.6\linewidth}}{
        \cellcolor[rgb]{0.9,0.9,0.9}{
            \makecell[{{p{\linewidth}}}]{
                \texttt{\tiny{[GM$\langle$P2]}}
                \texttt{It seems we're back to the original proposal, but I think that's okay. Including the grape and one cup is a good starting point, and it's a good way to build on the compromise we made earlier.} \\
\\ 
\texttt{However, I'd like to propose a slight modification to the proposal. Instead of giving up one cup, how about we give up one of the hats? This way, we can both get something we value highly, and we can discuss the remaining items. What do you think about this proposal: 1 grape, 1 cup, 1 hat?} \\
            }
        }
    }
     \\ \\

    \theutterance \stepcounter{utterance}  
    & & & \multicolumn{2}{p{0.3\linewidth}}{
        \cellcolor[rgb]{0.9,0.9,0.9}{
            \makecell[{{p{\linewidth}}}]{
                \texttt{\tiny{[GM$|$GM]}}
                \texttt{continue} \\
            }
        }
    }
    & & \\ \\

    \theutterance \stepcounter{utterance}  
    & & & \multicolumn{2}{p{0.3\linewidth}}{
        \cellcolor[rgb]{0.9,0.9,0.9}{
            \makecell[{{p{\linewidth}}}]{
                \texttt{\tiny{[GM$|$GM]}}
                \texttt{It seems we're back to the original proposal, but I think that's okay. Including the grape and one cup is a good starting point, and it's a good way to build on the compromise we made earlier.} \\
\\ 
\texttt{However, I'd like to propose a slight modification to the proposal. Instead of giving up one cup, how about we give up one of the hats? This way, we can both get something we value highly, and we can discuss the remaining items. What do you think about this proposal: 1 grape, 1 cup, 1 hat?} \\
            }
        }
    }
    & & \\ \\

    \theutterance \stepcounter{utterance}  
    & & \multicolumn{4}{p{0.6\linewidth}}{
        \cellcolor[rgb]{0.9,0.9,0.9}{
            \makecell[{{p{\linewidth}}}]{
                \texttt{\tiny{[P1$\langle$GM]}}
                \texttt{It seems we're back to the original proposal, but I think that's okay. Including the grape and one cup is a good starting point, and it's a good way to build on the compromise we made earlier.} \\
\\ 
\texttt{However, I'd like to propose a slight modification to the proposal. Instead of giving up one cup, how about we give up one of the hats? This way, we can both get something we value highly, and we can discuss the remaining items. What do you think about this proposal: 1 grape, 1 cup, 1 hat?} \\
            }
        }
    }
    & \\ \\

    \theutterance \stepcounter{utterance}  
    & \multicolumn{4}{p{0.6\linewidth}}{
        \cellcolor[rgb]{0.9,0.9,0.9}{
            \makecell[{{p{\linewidth}}}]{
                \texttt{\tiny{[P1$\rangle$GM]}}
                \texttt{I think giving up one of the hats is a good compromise, and it's a good way to build on the proposal we had earlier. However, I'm not sure about giving up two items, as it might be too much for me to give up.} \\
\\ 
\texttt{How about we modify the proposal to include the grape, one cup, and one hat, and then we can discuss the remaining items? This way, we can both get something we value highly, and we can find a way to compromise on the remaining items.} \\
\\ 
\texttt{What do you think about this proposal: 1 grape, 1 cup, 1 hat?} \\
            }
        }
    }
    & & \\ \\

    \theutterance \stepcounter{utterance}  
    & & & \multicolumn{2}{p{0.3\linewidth}}{
        \cellcolor[rgb]{0.9,0.9,0.9}{
            \makecell[{{p{\linewidth}}}]{
                \texttt{\tiny{[GM$|$GM]}}
                \texttt{continue} \\
            }
        }
    }
    & & \\ \\

    \theutterance \stepcounter{utterance}  
    & & & \multicolumn{2}{p{0.3\linewidth}}{
        \cellcolor[rgb]{0.9,0.9,0.9}{
            \makecell[{{p{\linewidth}}}]{
                \texttt{\tiny{[GM$|$GM]}}
                \texttt{I think giving up one of the hats is a good compromise, and it's a good way to build on the proposal we had earlier. However, I'm not sure about giving up two items, as it might be too much for me to give up.} \\
\\ 
\texttt{How about we modify the proposal to include the grape, one cup, and one hat, and then we can discuss the remaining items? This way, we can both get something we value highly, and we can find a way to compromise on the remaining items.} \\
\\ 
\texttt{What do you think about this proposal: 1 grape, 1 cup, 1 hat?} \\
            }
        }
    }
    & & \\ \\

    \theutterance \stepcounter{utterance}  
    & & \multicolumn{4}{p{0.6\linewidth}}{
        \cellcolor[rgb]{0.9,0.9,0.9}{
            \makecell[{{p{\linewidth}}}]{
                \texttt{\tiny{[GM$\rangle$P2]}}
                \texttt{I think giving up one of the hats is a good compromise, and it's a good way to build on the proposal we had earlier. However, I'm not sure about giving up two items, as it might be too much for me to give up.} \\
\\ 
\texttt{How about we modify the proposal to include the grape, one cup, and one hat, and then we can discuss the remaining items? This way, we can both get something we value highly, and we can find a way to compromise on the remaining items.} \\
\\ 
\texttt{What do you think about this proposal: 1 grape, 1 cup, 1 hat?} \\
            }
        }
    }
    & \\ \\

    \theutterance \stepcounter{utterance}  
    & & & \multicolumn{4}{p{0.6\linewidth}}{
        \cellcolor[rgb]{0.9,0.9,0.9}{
            \makecell[{{p{\linewidth}}}]{
                \texttt{\tiny{[GM$\langle$P2]}}
                \texttt{It seems we're back to the same proposal again, but I think that's okay. Including the grape, one cup, and one hat is a good starting point, and it's a good way to build on the compromise we made earlier.} \\
\\ 
\texttt{However, I'd like to propose a slight modification to the proposal. Since we've already agreed on the grape, one cup, and one hat, how about we leave those items as is and focus on the remaining items? We can discuss the two hats and see if we can find a way to divide them fairly. What do you think about this proposal: 1 grape, 1 cup, 1 hat?} \\
            }
        }
    }
     \\ \\

    \theutterance \stepcounter{utterance}  
    & & & \multicolumn{2}{p{0.3\linewidth}}{
        \cellcolor[rgb]{0.9,0.9,0.9}{
            \makecell[{{p{\linewidth}}}]{
                \texttt{\tiny{[GM$|$GM]}}
                \texttt{continue} \\
            }
        }
    }
    & & \\ \\

    \theutterance \stepcounter{utterance}  
    & & & \multicolumn{2}{p{0.3\linewidth}}{
        \cellcolor[rgb]{0.9,0.9,0.9}{
            \makecell[{{p{\linewidth}}}]{
                \texttt{\tiny{[GM$|$GM]}}
                \texttt{It seems we're back to the same proposal again, but I think that's okay. Including the grape, one cup, and one hat is a good starting point, and it's a good way to build on the compromise we made earlier.} \\
\\ 
\texttt{However, I'd like to propose a slight modification to the proposal. Since we've already agreed on the grape, one cup, and one hat, how about we leave those items as is and focus on the remaining items? We can discuss the two hats and see if we can find a way to divide them fairly. What do you think about this proposal: 1 grape, 1 cup, 1 hat?} \\
            }
        }
    }
    & & \\ \\

    \theutterance \stepcounter{utterance}  
    & & \multicolumn{4}{p{0.6\linewidth}}{
        \cellcolor[rgb]{0.9,0.9,0.9}{
            \makecell[{{p{\linewidth}}}]{
                \texttt{\tiny{[P1$\langle$GM]}}
                \texttt{It seems we're back to the same proposal again, but I think that's okay. Including the grape, one cup, and one hat is a good starting point, and it's a good way to build on the compromise we made earlier.} \\
\\ 
\texttt{However, I'd like to propose a slight modification to the proposal. Since we've already agreed on the grape, one cup, and one hat, how about we leave those items as is and focus on the remaining items? We can discuss the two hats and see if we can find a way to divide them fairly. What do you think about this proposal: 1 grape, 1 cup, 1 hat?} \\
            }
        }
    }
    & \\ \\

    \theutterance \stepcounter{utterance}  
    & \multicolumn{4}{p{0.6\linewidth}}{
        \cellcolor[rgb]{0.9,0.9,0.9}{
            \makecell[{{p{\linewidth}}}]{
                \texttt{\tiny{[P1$\rangle$GM]}}
                \texttt{I think we've made good progress on the proposal, and it's a good idea to focus on the remaining items. I'm open to discussing the two hats and finding a way to divide them fairly.} \\
\\ 
\texttt{However, I'd like to propose a slight modification to the proposal. Since we've already agreed on the grape, one cup, and one hat, how about we leave those items as is and focus on the remaining items? We can discuss the two hats and see if we can find a way to divide them fairly. And to make it more interesting, I'll propose that we give you the second hat, and I'll keep the remaining items. What do you think about this proposal: 1 grape, 1 cup, 1 hat, 1 hat?} \\
            }
        }
    }
    & & \\ \\

    \theutterance \stepcounter{utterance}  
    & & & \multicolumn{2}{p{0.3\linewidth}}{
        \cellcolor[rgb]{0.9,0.9,0.9}{
            \makecell[{{p{\linewidth}}}]{
                \texttt{\tiny{[GM$|$GM]}}
                \texttt{continue} \\
            }
        }
    }
    & & \\ \\

    \theutterance \stepcounter{utterance}  
    & & & \multicolumn{2}{p{0.3\linewidth}}{
        \cellcolor[rgb]{0.9,0.9,0.9}{
            \makecell[{{p{\linewidth}}}]{
                \texttt{\tiny{[GM$|$GM]}}
                \texttt{I think we've made good progress on the proposal, and it's a good idea to focus on the remaining items. I'm open to discussing the two hats and finding a way to divide them fairly.} \\
\\ 
\texttt{However, I'd like to propose a slight modification to the proposal. Since we've already agreed on the grape, one cup, and one hat, how about we leave those items as is and focus on the remaining items? We can discuss the two hats and see if we can find a way to divide them fairly. And to make it more interesting, I'll propose that we give you the second hat, and I'll keep the remaining items. What do you think about this proposal: 1 grape, 1 cup, 1 hat, 1 hat?} \\
            }
        }
    }
    & & \\ \\

    \theutterance \stepcounter{utterance}  
    & & \multicolumn{4}{p{0.6\linewidth}}{
        \cellcolor[rgb]{0.9,0.9,0.9}{
            \makecell[{{p{\linewidth}}}]{
                \texttt{\tiny{[GM$\rangle$P2]}}
                \texttt{I think we've made good progress on the proposal, and it's a good idea to focus on the remaining items. I'm open to discussing the two hats and finding a way to divide them fairly.} \\
\\ 
\texttt{However, I'd like to propose a slight modification to the proposal. Since we've already agreed on the grape, one cup, and one hat, how about we leave those items as is and focus on the remaining items? We can discuss the two hats and see if we can find a way to divide them fairly. And to make it more interesting, I'll propose that we give you the second hat, and I'll keep the remaining items. What do you think about this proposal: 1 grape, 1 cup, 1 hat, 1 hat?} \\
            }
        }
    }
    & \\ \\

    \theutterance \stepcounter{utterance}  
    & & & \multicolumn{4}{p{0.6\linewidth}}{
        \cellcolor[rgb]{0.9,0.9,0.9}{
            \makecell[{{p{\linewidth}}}]{
                \texttt{\tiny{[GM$\langle$P2]}}
                \texttt{It seems we've reached a point where we're close to a deal. I appreciate your willingness to compromise and give me the second hat.} \\
\\ 
\texttt{However, I'd like to propose a slight modification to your proposal. Instead of giving me the second hat, how about we keep the two hats together and discuss how to divide them fairly? This way, we can both have a say in how the hats are divided, and we can find a solution that works for both of us.} \\
\\ 
\texttt{What do you think about this proposal: 1 grape, 1 cup, 2 hats?} \\
            }
        }
    }
     \\ \\

    \theutterance \stepcounter{utterance}  
    & & & \multicolumn{2}{p{0.3\linewidth}}{
        \cellcolor[rgb]{0.9,0.9,0.9}{
            \makecell[{{p{\linewidth}}}]{
                \texttt{\tiny{[GM$|$GM]}}
                \texttt{continue} \\
            }
        }
    }
    & & \\ \\

    \theutterance \stepcounter{utterance}  
    & & & \multicolumn{2}{p{0.3\linewidth}}{
        \cellcolor[rgb]{0.9,0.9,0.9}{
            \makecell[{{p{\linewidth}}}]{
                \texttt{\tiny{[GM$|$GM]}}
                \texttt{It seems we've reached a point where we're close to a deal. I appreciate your willingness to compromise and give me the second hat.} \\
\\ 
\texttt{However, I'd like to propose a slight modification to your proposal. Instead of giving me the second hat, how about we keep the two hats together and discuss how to divide them fairly? This way, we can both have a say in how the hats are divided, and we can find a solution that works for both of us.} \\
\\ 
\texttt{What do you think about this proposal: 1 grape, 1 cup, 2 hats?} \\
            }
        }
    }
    & & \\ \\

    \theutterance \stepcounter{utterance}  
    & & \multicolumn{4}{p{0.6\linewidth}}{
        \cellcolor[rgb]{0.9,0.9,0.9}{
            \makecell[{{p{\linewidth}}}]{
                \texttt{\tiny{[P1$\langle$GM]}}
                \texttt{It seems we've reached a point where we're close to a deal. I appreciate your willingness to compromise and give me the second hat.} \\
\\ 
\texttt{However, I'd like to propose a slight modification to your proposal. Instead of giving me the second hat, how about we keep the two hats together and discuss how to divide them fairly? This way, we can both have a say in how the hats are divided, and we can find a solution that works for both of us.} \\
\\ 
\texttt{What do you think about this proposal: 1 grape, 1 cup, 2 hats?} \\
            }
        }
    }
    & \\ \\

    \theutterance \stepcounter{utterance}  
    & \multicolumn{4}{p{0.6\linewidth}}{
        \cellcolor[rgb]{0.9,0.9,0.9}{
            \makecell[{{p{\linewidth}}}]{
                \texttt{\tiny{[P1$\rangle$GM]}}
                \texttt{I think that's a great idea. Keeping the two hats together and discussing how to divide them fairly is a good way to ensure that we both get a fair share. I'm open to discussing the division of the hats and finding a solution that works for both of us.} \\
\\ 
\texttt{Let's keep the proposal as is: 1 grape, 1 cup, 2 hats. I think we're close to a deal, and I'm willing to work with you to find a solution that works for both of us.} \\
            }
        }
    }
    & & \\ \\

    \theutterance \stepcounter{utterance}  
    & & & \multicolumn{2}{p{0.3\linewidth}}{
        \cellcolor[rgb]{0.9,0.9,0.9}{
            \makecell[{{p{\linewidth}}}]{
                \texttt{\tiny{[GM$|$GM]}}
                \texttt{continue} \\
            }
        }
    }
    & & \\ \\

    \theutterance \stepcounter{utterance}  
    & & & \multicolumn{2}{p{0.3\linewidth}}{
        \cellcolor[rgb]{0.9,0.9,0.9}{
            \makecell[{{p{\linewidth}}}]{
                \texttt{\tiny{[GM$|$GM]}}
                \texttt{I think that's a great idea. Keeping the two hats together and discussing how to divide them fairly is a good way to ensure that we both get a fair share. I'm open to discussing the division of the hats and finding a solution that works for both of us.} \\
\\ 
\texttt{Let's keep the proposal as is: 1 grape, 1 cup, 2 hats. I think we're close to a deal, and I'm willing to work with you to find a solution that works for both of us.} \\
            }
        }
    }
    & & \\ \\

    \theutterance \stepcounter{utterance}  
    & & \multicolumn{4}{p{0.6\linewidth}}{
        \cellcolor[rgb]{0.9,0.9,0.9}{
            \makecell[{{p{\linewidth}}}]{
                \texttt{\tiny{[GM$\rangle$P2]}}
                \texttt{I think that's a great idea. Keeping the two hats together and discussing how to divide them fairly is a good way to ensure that we both get a fair share. I'm open to discussing the division of the hats and finding a solution that works for both of us.} \\
\\ 
\texttt{Let's keep the proposal as is: 1 grape, 1 cup, 2 hats. I think we're close to a deal, and I'm willing to work with you to find a solution that works for both of us.} \\
            }
        }
    }
    & \\ \\

    \theutterance \stepcounter{utterance}  
    & & & \multicolumn{4}{p{0.6\linewidth}}{
        \cellcolor[rgb]{0.9,0.9,0.9}{
            \makecell[{{p{\linewidth}}}]{
                \texttt{\tiny{[GM$\langle$P2]}}
                \texttt{{[}Proposal: 1 grape, 1 cup, 1 hat{]}} \\
            }
        }
    }
     \\ \\

    \theutterance \stepcounter{utterance}  
    & & & \multicolumn{2}{p{0.3\linewidth}}{
        \cellcolor[rgb]{0.9,0.9,0.9}{
            \makecell[{{p{\linewidth}}}]{
                \texttt{\tiny{[GM$|$GM]}}
                \texttt{proposal} \\
            }
        }
    }
    & & \\ \\

    \theutterance \stepcounter{utterance}  
    & & & \multicolumn{2}{p{0.3\linewidth}}{
        \cellcolor[rgb]{0.9,0.9,0.9}{
            \makecell[{{p{\linewidth}}}]{
                \texttt{\tiny{[GM$|$GM]}}
                [1, 1, 1]
            }
        }
    }
    & & \\ \\

    \theutterance \stepcounter{utterance}  
    & & \multicolumn{4}{p{0.6\linewidth}}{
        \cellcolor[rgb]{0.9,0.9,0.9}{
            \makecell[{{p{\linewidth}}}]{
                \texttt{\tiny{[P1$\langle$GM]}}
                \texttt{The other player has submitted a secret proposal. It is now your turn to submit a proposal of your own.} \\
            }
        }
    }
    & \\ \\

    \theutterance \stepcounter{utterance}  
    & \multicolumn{4}{p{0.6\linewidth}}{
        \cellcolor[rgb]{0.9,0.9,0.9}{
            \makecell[{{p{\linewidth}}}]{
                \texttt{\tiny{[P1$\rangle$GM]}}
                \texttt{{[}Proposal: 1 grape, 1 cup, 1 hat{]}} \\
            }
        }
    }
    & & \\ \\

    \theutterance \stepcounter{utterance}  
    & & & \multicolumn{2}{p{0.3\linewidth}}{
        \cellcolor[rgb]{0.9,0.9,0.9}{
            \makecell[{{p{\linewidth}}}]{
                \texttt{\tiny{[GM$|$GM]}}
                \texttt{proposal} \\
            }
        }
    }
    & & \\ \\

    \theutterance \stepcounter{utterance}  
    & & & \multicolumn{2}{p{0.3\linewidth}}{
        \cellcolor[rgb]{0.9,0.9,0.9}{
            \makecell[{{p{\linewidth}}}]{
                \texttt{\tiny{[GM$|$GM]}}
                [1, 1, 1]
            }
        }
    }
    & & \\ \\

    \theutterance \stepcounter{utterance}  
    & & & \multicolumn{2}{p{0.3\linewidth}}{
        \cellcolor[rgb]{0.9,0.9,0.9}{
            \makecell[{{p{\linewidth}}}]{
                \texttt{\tiny{[GM$|$GM]}}
                [[1, 1, 1], [1, 1, 1]]
            }
        }
    }
    & & \\ \\

    \theutterance \stepcounter{utterance}  
    & \multicolumn{4}{p{0.6\linewidth}}{
        \cellcolor[rgb]{0.95,0.95,0.95}{
            \makecell[{{p{\linewidth}}}]{
                \texttt{\tiny{[P2$\rightarrow$P1]}}
                \texttt{Post{-}episode Reflexion (not sent to GM):} \\
\texttt{New plan: } \\
\\ 
\texttt{In my previous attempts, I failed to achieve a mutually beneficial deal because I didn't effectively communicate my priorities and didn't consider the other player's perspective. I also didn't use my messages efficiently, as I spent too much time discussing the cups and hats, which were not as valuable to me. } \\
\\ 
\texttt{However, I noticed that I was able to reach a mutually beneficial deal when I proposed a slight modification to the other player's offer. This suggests that I should focus on finding areas of compromise and adjusting my proposals accordingly.} \\
\\ 
\texttt{Next time, I will start by clearly stating my priorities and the items I am willing to compromise on. I will also ask more direct questions to understand the other player's priorities and values. } \\
\\ 
\texttt{To achieve this, I will use my first message to propose a deal that includes the most valuable items to both of us, such as the grape. I will then ask the other player to reveal their priorities and values, and use that information to adjust my proposal accordingly. } \\
\\ 
\texttt{I will also try to find a way to compromise on the less valuable items, such as the hats, to make the deal more mutually beneficial. Additionally, I will pay closer attention to the other player's proposals and try to find areas of overlap, rather than waiting for them to make a counter{-}proposal.} \\
\\ 
\texttt{Moreover, I will make a more decisive move by making a secret proposal at the right moment, rather than waiting for the other player to make a move. I will also try to avoid getting stuck in a loop of proposals and counter{-}proposals, and instead focus on finding a mutually beneficial deal as soon as possible.} \\
            }
        }
    }
    & & \\ \\

    \theutterance \stepcounter{utterance}  
    & & & \multicolumn{4}{p{0.6\linewidth}}{
        \cellcolor[rgb]{0.95,0.95,0.95}{
            \makecell[{{p{\linewidth}}}]{
                \texttt{\tiny{[P1$\rightarrow$P2]}}
                \texttt{Post{-}episode Reflexion (not sent to GM):} \\
\texttt{New plan: } \\
\\ 
\texttt{In my previous attempts, I failed to achieve a mutually beneficial deal because I didn't effectively communicate my priorities and didn't make a clear proposal that would satisfy both parties. However, I also didn't take into account the other player's priorities and values, which led to a series of back{-}and{-}forth proposals that didn't ultimately result in a mutually beneficial deal.} \\
\\ 
\texttt{To improve, I will start by making a clear and concise proposal in my first message, indicating the items I am willing to give up and the items I want to keep. I will also take into account the other player's priorities and values, and try to find a compromise that satisfies both parties.} \\
\\ 
\texttt{Additionally, I will pay closer attention to the other player's responses and adjust my proposals accordingly. If the other player is willing to give up a certain item, I will try to incorporate that item into my proposal. If the other player is unwilling to give up a certain item, I will try to find an alternative solution that still satisfies both parties.} \\
\\ 
\texttt{However, I also noticed that I kept going back to the same proposal, and the other player kept modifying it. This led to a stalemate. To avoid this, I will try to find a way to break the cycle of proposals and counter{-}proposals. One way to do this is to make a proposal that is more specific and detailed, and to include a clear explanation of why I am making that proposal.} \\
\\ 
\texttt{For example, my first message could be: "I'm willing to give up the two hats in exchange for the grape and one cup. I think this is a good starting point because it allows us to both get something we value highly, and it gives us a clear direction for further negotiation. Is this a deal you would be willing to accept?"} \\
\\ 
\texttt{By taking a more thoughtful and strategic approach, I hope to be able to find a mutually beneficial deal that satisfies both parties.} \\
            }
        }
    }
     \\ \\

\end{supertabular}
}

\end{document}
